\section{Introduction}
\label{sec:intro}

A benchmark number for a large language model (LLM) is the product of the model \emph{and} the decoder that turns its next-token distributions into text. Decoding settings are rarely held fixed across reports. Leaderboards use greedy decoding, chat deployments use temperatures near 0.7, and code papers report pass@1 sampled at temperatures between 0.2 and 0.8~\citep{chen2021codex}. Many papers do not state the setting at all. If a domain's scores depend strongly on the decoder, cross-paper comparisons in that domain are unreliable. If its \emph{model rankings} depend on the decoder, even the qualitative conclusion ``model A beats model B'' is in question.

{\bf Why does this matter now?} Practitioners routinely choose between greedy decoding and a range of samplers (temperature, top-$k$, nucleus/top-$p$~\citep{holtzman2020curious}, and the more recent min-$p$~\citep{nguyen2024minp}) to trade accuracy for diversity. They need to know which tasks tolerate sampling and which sampler loses the least. Evaluators need to know whether a ranking computed under one decoder carries over to another.

{\bf What is missing?} Prior work mostly studies decoding one task family at a time. Nucleus sampling was proposed for open-ended generation~\citep{holtzman2020curious}, temperature was studied against pass@$k$ for code~\citep{chen2021codex}, and temperature was found to have little effect on multiple-choice problem solving~\citep{renze2024temperature}. The broadest recent examination~\citep{shi2024thorough} covers many tasks but few models and does not analyze ranking stability. To our knowledge, no single controlled grid of models $\times$ domains $\times$ samplers $\times$ seeds measures score sensitivity, ranking stability, and seed variance side by side.

{\bf Our approach.} We run 11 open-weight instruction-tuned LLMs from 7 families on four domains: factual QA (\triviaqa), multi-step math reasoning (\gsm), code synthesis (\humaneval), and creative story writing (\rocstories prompts scored by an LLM judge). Each model is decoded with six strategies, and every sampled strategy is run with three seeds. This gives 132{,}064 generations, all served locally with \vllm~\citep{kwon2023vllm}. We pre-register five quantitative hypotheses (H1--H5 in \secref{sec:method}) and test each one with paired statistics and item-bootstrap confidence intervals.

{\bf What do we find?} Sampling never helps on average in a verifiable domain. Pure \tone sampling costs 4.3 points on factual QA (an 11.0\% relative drop), 5.9 on \gsm, and 6.7 on \humaneval. \topk behaves almost exactly like \tone. Truncated samplers are much safer: \topp recovers 50\% of the \gsm loss and \minp recovers 57\%. Factual QA is \emph{not} decoding-invariant; its relative loss matches that of reasoning. Sampling does not raise LLM-judged creative quality: changes range from $+0.0\%$ to $+0.7\%$ and are not significant, and \tone costs $7.9\%$. Model rankings are reassuringly stable. \ktau against the greedy ranking is 0.93 on factual QA, 0.85--0.89 on reasoning, and 0.93--1.00 on code. The only $\tau$ near 0.7 comes from small models degenerating on creative writing at \tone. Code is the noisiest domain across seeds, with a seed standard deviation of 2.13 points and 26.7\% of problems flipping outcome. About half of that gap disappears once we control for benchmark size.

Our contributions are:
\begin{itemize}[leftmargin=*,itemsep=0pt,topsep=0pt]
    \item We conduct a controlled 11-model $\times$ 4-domain $\times$ 6-strategy $\times$ 3-seed study of decoding sensitivity, with all samplers implemented identically in one serving stack.
    \item We test five quantitative hypotheses about domain-specific sensitivity and refute or qualify four of them. Factual QA is not robust to sampling (H1), the reasoning gap is strongly model-dependent (H2), sampling does not improve judged creativity (H3), and reasoning rankings are less stable than factual ones but far from $\tau<0.7$ (H4).
    \item We measure ranking stability with item-bootstrap \ktau and seed variability with a matched-$n$ control. We show that rankings shift mainly when models are tightly packed or when small models degenerate.
    \item We derive practical recommendations: use greedy or min-$p$/top-$p$ for evaluation, never raw \tone or \topk, and average at least three seeds when sampling, especially for code.
\end{itemize}

The rest of the paper is organized as follows. \Secref{sec:related} reviews related work, \secref{sec:method} describes the experimental design, \secref{sec:results} presents results, \secref{sec:discussion} discusses implications and limitations, and \secref{sec:conclusion} concludes.
